% Einheit 4 von 4 – Terme und Sachaufgaben (bank/rationale-zahlen/e4.jsonl, zusammenbau v0.1)
%% TODO Zweigzeile Teil 1: was hier gelernt wird (Satz in Schülersprache aus dem Einheitstitel) – Einheit 4
\einheitenkopf[e4]{Einheit 4 von 4 · Terme und Sachaufgaben}
\zweigzeile{ab Kl. 7 · P10}
\setcounter{aufgabe}{22}

%% TODO Titel: Ich-kann-Satz fehlt in der Bank (Platzhalter: Kettenname) – Nr. 23
%% TODO Anweisung: ein Satz über der ersten Teilaufgabe fehlt in der Bank – Nr. 23
\begin{aufgabe}{Terme und Sachaufgaben}
%% TODO \rechenplatz{n}: Schrittzahl des Grundfalls fehlt in der Bank (nur bei fester Schreibform, 2.3 b)
\begin{teile}
\teil Auf Mias Konto sind $120$ €. Dann werden $45$ € abgebucht. Was ist der Startwert, was die Änderung? Startwert \leerfeld , Änderung \leerfeld
\teil Kontostand $-35$ €, Einzahlung $50$ € – neuer Kontostand? \leerfeld[€]
\teil Kontostand $60$ €; Buchungen $-27$ €, $-48$ €, $+30$ € – neuer Kontostand? \leerfeld[€]
\teil Morgens $-9$ °C, mittags $4$ °C – um wie viel Grad wärmer? \leerfeld[Grad]
\teil Löse die Klammer auf und berechne: $15 - (8 - 17)$ \leerfeld
\teil $3 \cdot (y + 4)$ für $y = -7$ – welcher Wert? \leerfeld
\teil Tierpark, Preise: Erwachsene $17$ €; Kinder von 4 bis 15 Jahren $8$ €; Kinder unter 4 Jahren frei; Familienkarte (2 Erwachsene und bis zu 3 Kinder) $46$ €; Großeltern-Enkel-Karte (2 Großeltern mit 1 oder 2 Enkeln, nur Montag bis Freitag) $30$ €. Am Mittwoch kommen Mutter, Vater, Oma, Opa und die Kinder Lina (3 Jahre), Paul (8) und Emma (9). Ermittle den günstigsten Gesamtpreis und vergleiche mit den anderen Möglichkeiten. \hfill (P10 2016 OS) \\ Günstigster Preis: \leerfeld[€]
\teil 2024 kamen $48\,315$ Menschen zum Stadtfest, das sind $3\,870$ mehr als 2023. Wie viele waren es 2023? \hfill (P10 2017 OS) \\ \leerfeld Menschen
\end{teile}
\end{aufgabe}

%% TODO Titel: Ich-kann-Satz fehlt in der Bank (Platzhalter: Kettenname) – Nr. 24
%% TODO Anweisung: ein Satz über der ersten Teilaufgabe fehlt in der Bank – Nr. 24
\begin{aufgabe}{Rechenvorteile (Tauschen, geschickt zusammenfassen)}
\begin{teile}
\teil Rechne geschickt: $-17 + 38 + 17$ \leerfeld
\end{teile}
\end{aufgabe}

%% TODO Titel: Ich-kann-Satz fehlt in der Bank (Platzhalter: Kettenname) – Nr. 25
%% TODO Anweisung: ein Satz über der ersten Teilaufgabe fehlt in der Bank – Nr. 25
\begin{aufgabe}{Terme und Sachaufgaben – Fehler finden, Begründen, Anwendung}
\begin{teile}
\teil Jan rechnet: \rechnung{30 - (6 - 9) &= 30 - 6 - 9 \\ &= 15} Finde den Fehler.
\teil Erkläre, warum man bei $-38 + 17 + 38$ zuerst $-38$ und $38$ zusammenfasst.
\teil Lea hat $85$ € auf dem Konto. Sie zahlt $24{,}99$ € Handyrechnung und $49{,}90$ € für ein Konzertticket, bekommt $30$ € Taschengeld und zahlt $11{,}50$ € fürs Kino. Überziehen darf sie nicht. Kann sie danach noch ein Buch für $15$ € kaufen?
\end{teile}
\end{aufgabe}
